% Problem-section file following ../_template.tex.
\section{Quantum capacity of the Gaussian random-displacement channel}
\label{sec:problem-64}

\paragraph{Problem.}
What is the unconstrained, unassisted quantum capacity of the single-mode
Gaussian random-displacement channel?  For a noise standard deviation
$\sigma>0$, define the channel by
\begin{equation}
  \mathcal N_\sigma(\rho)
  :=\frac{1}{\pi\sigma^2}\int_{\mathbb C}
       e^{-\lvert\alpha\rvert^2/\sigma^2}
       D(\alpha)\rho D(\alpha)^\dagger\,d^2\alpha,
  \qquad
  D(\alpha):=e^{\alpha a^\dagger-\alpha^*a},
  \label{eq:p64-random-displacement-channel}
\end{equation}
where $a$ is the mode annihilation operator.  Equivalently, writing
$\alpha=(x+iy)/\sqrt2$, the channel in
Eq.~\eqref{eq:p64-random-displacement-channel} adds independent classical
Gaussian shifts $x,y\sim\mathcal N(0,\sigma^2)$ to the two canonical
quadratures.  With $\hat n_j:=a_j^\dagger a_j$, define its
energy-unconstrained capacity as
\begin{equation}
  \begin{aligned}
  \mathcal Q(\mathcal N_\sigma)
  &:={\sup}_{0<N_{\rm S}<\infty}\,
     \lim_{n\to\infty}\frac1n
     {\sup}_{\substack{\rho_{A^n}:\\
       \operatorname{Tr}[\rho_{A^n}\sum_{j=1}^n\hat n_j]
       \leq nN_{\rm S}}}
       I_{\rm c}(\rho_{A^n},\mathcal N_\sigma^{\otimes n}),\\
  I_{\rm c}(\rho,\mathcal M)
  &:=S(\mathcal M(\rho))-S(\mathcal M^{\rm c}(\rho)),
  \end{aligned}
  \label{eq:p64-quantum-capacity}
\end{equation}
where $\mathcal M^{\rm c}$ is a complementary channel and logarithms in the
entropy are base two.  Determine Eq.~\eqref{eq:p64-quantum-capacity} for every
$\sigma>0$.

\subsection*{Status}

Unsolved

\subsection*{Source}

Lin and Noh explicitly describe determining
$\mathcal Q(\mathcal N_\sigma)$ as a longstanding open problem
\sourcecite{ref:p64-lin-noh}{LN25}.  Subramanian, Zheng, and Jiang subsequently
reaffirm that the exact capacity remains unknown while constructing codes
that attain the standard coherent-information rate
\sourcecite{ref:p64-subramanian-zheng-jiang}{SZJ25}.

\subsection*{Progress}

\begin{itemize}
  \item In the convention of
  Eq.~\eqref{eq:p64-random-displacement-channel}, the standard analytic bounds
  are
  \begin{equation}
    \left[\log_2\!\frac{1}{e\sigma^2}\right]_+
    \leq \mathcal Q(\mathcal N_\sigma)
    \leq U(\sigma),
    \qquad
    U(\sigma):=
    \begin{cases}
      \displaystyle\log_2\!\frac{1-\sigma^2}{\sigma^2},
        &0<\sigma<1/\sqrt2,\\[1mm]
      0,&\sigma\geq1/\sqrt2,
    \end{cases}
    \qquad [t]_+:=\max\{t,0\}.
    \label{eq:p64-standard-bounds}
  \end{equation}
  The lower bound in Eq.~\eqref{eq:p64-standard-bounds} is the infinite-energy
  limit of one-use coherent information, while the upper bound follows from a
  channel decomposition and data processing
  \sourcecite{ref:p64-holevo-werner}{HW01},
  \sourcecite{ref:p64-noh-albert-jiang}{NAJ19}.  Thus the capacity is positive
  for $0<\sigma<1/\sqrt e$ and vanishes for
  $\sigma\geq1/\sqrt2$, but the bounds do not determine it in the remaining
  positive-capacity regime; for
  $1/\sqrt e\leq\sigma<1/\sqrt2$, even positivity is unknown.  Degradable
  flagged extensions sharpen the upper bound at low noise without closing the
  gap \sourcecite{ref:p64-fanizza-kianvash-giovannetti}{FKG21}.

  \item Lin and Noh evaluated multimode GKP families with maximum-likelihood
  decoding.  Their largest surface--square instance, of distance $39$, has a
  positive computed hashing rate at $\sigma=0.598$, and its finite-size
  crossings indicate a threshold close to $1/\sqrt e$.  These are numerical
  code-performance results below the known positivity threshold, not an exact
  capacity evaluation or a proof of positivity beyond $1/\sqrt e$
  \sourcecite{ref:p64-lin-noh}{LN25}.

  \item Subramanian, Zheng, and Jiang constructed concatenated square-GKP and
  quantum polar codes with analog decoding whose asymptotic rate is
  \begin{equation}
    R_{\rm GKP\text{-}polar}(\sigma)
    =\left[\log_2\!\frac{1}{e\sigma^2}\right]_+.
    \label{eq:p64-gkp-polar-rate}
  \end{equation}
  Equation~\eqref{eq:p64-gkp-polar-rate} removes earlier discreteness and
  nonconstructiveness restrictions in attaining the lower bound of
  Eq.~\eqref{eq:p64-standard-bounds}.  It does not prove that this lower bound
  equals the regularized capacity in Eq.~\eqref{eq:p64-quantum-capacity}
  \sourcecite{ref:p64-subramanian-zheng-jiang}{SZJ25}.
\end{itemize}

\subsection*{References}

\begin{enumerate}[label={},leftmargin=4.8em,labelsep=0.5em]
  \footnotesize
  \item[\textup{[HW01]}]\label{ref:p64-holevo-werner}
  A. S. Holevo and R. F. Werner,
  ``Evaluating Capacities of Bosonic Gaussian Channels,''
  \emph{Physical Review A} \textbf{63}, 032312 (2001).
  \href{https://doi.org/10.1103/PhysRevA.63.032312}%
       {doi:10.1103/PhysRevA.63.032312};
  \href{https://arxiv.org/abs/quant-ph/9912067}%
       {arXiv:quant-ph/9912067}.

  \item[\textup{[NAJ19]}]\label{ref:p64-noh-albert-jiang}
  {\sloppy
  K. Noh, V. V. Albert, and L. Jiang,
  ``Quantum Capacity Bounds of Gaussian Thermal Loss Channels and Achievable
  Rates with Gottesman--Kitaev--Preskill Codes,''
  \emph{IEEE Transactions on Information Theory} \textbf{65}, 2563--2582
  (2019).
  \href{https://doi.org/10.1109/TIT.2018.2873764}%
       {doi:10.1109/TIT.2018.2873764};
  \href{https://arxiv.org/abs/1801.07271}{arXiv:1801.07271}.
  \par}

  \item[\textup{[FKG21]}]\label{ref:p64-fanizza-kianvash-giovannetti}
  M. Fanizza, F. Kianvash, and V. Giovannetti,
  ``Estimating Quantum and Private Capacities of Gaussian Channels via
  Degradable Extensions,'' \emph{Physical Review Letters} \textbf{127},
  210501 (2021).
  \href{https://doi.org/10.1103/PhysRevLett.127.210501}%
       {doi:10.1103/PhysRevLett.127.210501};
  \href{https://arxiv.org/abs/2103.09569}{arXiv:2103.09569}.

  \item[\textup{[LN25]}]\label{ref:p64-lin-noh}
  M. Lin and K. Noh,
  ``Exploring the Quantum Capacity of a Gaussian Random-Displacement Channel
  Using Gottesman--Kitaev--Preskill Codes and Maximum-Likelihood Decoding,''
  \emph{Physical Review A} \textbf{111}, 052445 (2025).
  \href{https://doi.org/10.1103/PhysRevA.111.052445}%
       {doi:10.1103/PhysRevA.111.052445};
  \href{https://arxiv.org/abs/2411.04277}{arXiv:2411.04277}.

  \item[\textup{[SZJ25]}]\label{ref:p64-subramanian-zheng-jiang}
  M. Subramanian, G. Zheng, and L. Jiang,
  ``Achievable Rates for Concatenated Square Gottesman--Kitaev--Preskill
  Codes,'' \emph{PRX Quantum} \textbf{6}, 040342 (2025).
  \href{https://doi.org/10.1103/56vj-z7h1}{doi:10.1103/56vj-z7h1};
  \href{https://arxiv.org/abs/2505.10499}{arXiv:2505.10499}.
\end{enumerate}

\subsection*{Comment}

The word ``optimal'' in the coding result of
\sourcecite{ref:p64-subramanian-zheng-jiang}{SZJ25} refers to attaining the
one-use coherent-information lower bound for all noise strengths, not to
proving that this rate is the quantum capacity.  The exact value remains
unknown throughout $0<\sigma<1/\sqrt2$; in the subinterval
$1/\sqrt e\leq\sigma<1/\sqrt2$, the first unresolved question is whether the
capacity is nonzero at all.

\subsection*{Tag}

Bosonic channels; Bosonic codes; Quantum capacity; Coherent information;
Gaussian quantum information; Continuous-variable systems

\subsection*{ID}

\texttt{op\_1e7f431ed9013d1e}
